% Audited optional addition to C29; not included automatically during source freeze.
\paragraph{Coefficient extension and a sharper local range.}
The relative assertion of Statement~\ref{sa08:C29} holds over every
commutative ring $k$, with $P=K(k,n)$ and the same finite comparison
bounds. Over a $\Z_{(p)}$-algebra it holds in the larger range $n<pq$.
The corresponding absolute finite-arity extension ranges are
$n\leq2q$ over arbitrary $k$ and $n\leq pq$ over a
$\Z_{(p)}$-algebra.

Here is the base-change justification, which does not assume that $k$
is flat over $\Z$. For $R=\Z$ or $\Z_{(p)}$, let $Q_R$ be the
bounded universal arity quotient evaluated on $K(R,n)$, and use the
analogous notation for the relative quotient. Every term is degreewise
free over $R$. Evaluation of a projective simplicial-module resolution
is an $R$-projective resolution, so these flat terms have no higher
derived extension-of-scalars groups. The finite cochain filtration gives
$k\otimes_R Q_R\simeq k\otimes_R^LQ_R$. Free divided-power base
change identifies the ordinary tensor with $Q_k$. Restriction of
scalars is exact, and derived adjunction therefore gives
\[
 R\operatorname{Hom}_{s\operatorname{Mod}_k}(Q_k,\underline k)
 \simeq
 R\operatorname{Hom}_{s\operatorname{Mod}_R}(Q_R,\underline k).
\]
The same holds for the relative quotient. All projective-dimension
bounds in the proof of C29 hold against every target $R$-module,
including $k$, proving the arbitrary-base assertion.

Over $R=\Z_{(p)}$, a variable multiplicity not divisible by $p$ is
a unit in the Euler homotopy. Thus every surviving component uses
at most $\lfloor m/p\rfloor$ variables, rather than
$\lfloor m/2\rfloor$. Its normalized support is at most
$n\lfloor m/p\rfloor$ and its simplicial projective dimension is
at most $n\lfloor m/p\rfloor+1$. Nonzero vertical cohomology still
starts in degree $qm+1$ over the torsionfree PID $R$.
These inequalities give positive absolute derived-Hom vanishing
when $n\leq pq$, and degree-zero absolute vanishing, hence relative
extension, when $n<pq$. If $p\nmid m$, total Euler already contracts
the arity quotient. The derived base-change argument transfers these
conclusions to every $R$-algebra. The prescribed full boundary and
finite-arity weight hypotheses are unchanged.
